\begin{abstract}
Single-cell RNA sequencing enables dissection of transcriptional heterogeneity at
per-cell resolution, yet perturbation studies commonly rely on principal component
analysis and linear classifiers that collapse the typed relational structure among
cells, genes, and regulatory modules.
We present a Heterogeneous Graph Transformer (HGT) framework for perturbation-aware
analysis of 9,201 MCF10A single-cell expression profiles across six joint p53/CENP-A
perturbation conditions (ArrayExpress E-MTAB-9861; \citep{jeffery2022cenpa}).
Our pipeline constructs a three-node-type heterogeneous graph linking cells via
k-nearest-neighbor neighborhoods, genes via coexpression, and both to
expression-derived gene modules, then applies a multitask HGT to classify p53
status, CENP-A status, and joint six-condition perturbation state.
Five-fold stratified cross-validation yields condition macro-F1 of
$0.981 \pm 0.003$ (logistic regression), $0.978 \pm 0.005$ (MLP), and $0.923 \pm 0.006$
(random forest); the multitask HGT reaches 0.967 on the fixed held-out split.
Silhouette-score analysis of learned cell representations shows that HGT embeddings
achieve a score of 0.782 for the six perturbation conditions, 40-fold higher than
raw PCA coordinates (0.020), indicating that near-ceiling linear classification
accuracy on strongly structured perturbation data does not reflect equivalent
representation quality.
UMAP visualization \citep{mcinnes2018umap} and confusion matrix analysis confirm
that residual misclassification concentrates at the p53-WT/DN boundary within the
chronic CENP-A overexpression condition.
Label-efficiency experiments show that logistic regression reaches 0.947 macro-F1
with only 10\% of labeled training data, reflecting the large perturbation effect
sizes in this dataset.
All code, figures, and \LaTeX{} sources are fully reproducible from the raw H5AD file.
\end{abstract}
